% =====================================================================
%  Chapter 4  ·  Session 4 --- Networking and firewall controls
% =====================================================================
\chapter{Networking and firewall controls}

Session~3 measured how much traffic the service survives; this chapter
decides which traffic may reach it at all. It covers addressing and ports,
the software-defined networks of the cloud, the firewall and its variants,
and defence in depth and Zero Trust. In the laboratory the VM of Session~2,
now serving Nginx, receives a default-deny policy: SSH only from the
administrator, the web port for everyone, nothing else.

\section{Addresses, ports and the attack surface}

Every firewall rule is written in the vocabulary of addressing. Two
machines on a network find each other by \emph{IP address}, and on one
machine a service is reached at a \emph{port}, a sixteen-bit number that
identifies which program the traffic is for. A web server listens on
port~80 (HTTP) or 443 (HTTPS), SSH on 22, a database on 5432 or 3306. The
pair \emph{(address, port)} therefore names one service on one machine, and
the combination \emph{(protocol, address, port)} on both ends, carried over
TCP (connection-oriented, ordered) or UDP (connectionless), is what a
firewall rule matches. \texttt{ss -tulpn} lists one's own side of this:
every port the machine listens on.

\begin{definition}[Attack surface]\label{def:attacksurface}
The \textbf{attack surface} of a system describes the sum of the points at
which an attacker can try to interact with it: every open port, every
reachable service, every exposed endpoint. Reducing the attack surface means
removing such points, since a service that is not listening, or a port that is
not reachable, cannot be attacked through that route.
\end{definition}

Comer (2021) uses the term \emph{security attack surface} in the same sense
for microservices and remote access (Comer, 2021, pp.~167--168, 234). Every
listening port is a door, some doors must be open for the service to work,
and the security question is which. For this reason the first useful action
on any machine is to run \texttt{ss -tulpn} and to justify each line.

\section{Networks in the cloud are software}

In a data centre the network a VM sits on is defined in software (Comer,
2021, Chapter~7). Because many tenants share one physical fabric, each
tenant gets a \emph{virtual network} built as an \emph{overlay} on the
shared physical \emph{underlay} (Comer, 2021, pp.~88--89). VLANs, and their
data-centre-scale successor VXLAN, tag traffic so that one tenant's packets
never mix with another's (Comer, 2021, pp.~89--90). A \emph{virtual switch}
inside each physical server connects the VMs running on it (Comer, 2021,
p.~91). \emph{Network address translation} (NAT) lets privately addressed
machines reach the outside world (Comer, 2021, pp.~91--92).

Boundaries between networks are therefore logical and enforced by
configuration, with two consequences: they can be misconfigured, and they
can be audited by reading the configuration.

\begin{definition}[Software-defined networking]\label{def:sdn}
\textbf{Software-defined networking} (SDN) describes the separation of a
network's \emph{control} (the decisions about where traffic goes) from its
\emph{forwarding} (the act of moving packets), so that a program, the
\emph{controller}, configures and monitors the switches. In the cloud this is
why an entire network can be created, changed or destroyed through an API in
seconds, as a VM can.
\end{definition}

Comer (2021) describes OpenFlow, the first SDN protocol, and a second
generation in which switches run small programs of their own (Comer, 2021,
pp.~94--96). Because the network is software, security is enforced in the
fabric itself, per resource, rather than by a physical appliance the traffic
happens to pass (CSA, 2024, p.~154). The tool for that is the security
group.

\begin{definition}[Security group]\label{def:secgroup}
A \textbf{security group} describes a virtual firewall enforced by the cloud
network fabric at the level of an instance, network interface or subnet. Its
rules allow or deny traffic by IP address, port and protocol. It is
\emph{stateful}: if a rule permits an outgoing request, the corresponding reply
is allowed back automatically.
\end{definition}

Two details matter in practice. First, security groups are stateful
and attach to resources, whereas \emph{network access control lists} (NACLs)
work lower in the stack, are typically stateless and apply to subnets (CSA,
2024, p.~156); the distinction decides which tool fits which job. Secondly,
providers differ in their defaults. AWS starts from a deny-all position, so
the customer must add allow rules, and two resources in the same security
group cannot talk to each other until a rule permits it (CSA, 2024, p.~158).
Azure, in contrast, starts from a more permissive position: the built-in
default rules of a network security group allow traffic freely within the
virtual network (and from the Azure load balancer), while a final default
rule denies inbound traffic from the Internet. The Internet is thus closed
by default, but east--west traffic between the customer's own resources is
open unless restricted, and that is exactly the lateral movement of
Session~2. A provider's defaults must therefore be known before they are
trusted, because assuming the wrong one leaves a service reachable from an
unintended place.

\section{The firewall}

\begin{figure}[htbp]
\centering
\begin{tikzpicture}[font=\footnotesize, node distance=7mm]
  \node[font=\small\itshape, text=kgrey, align=center] (in) at (0,0) {incoming\\traffic};
  \node[redbox, right=10mm of in] (fw) {firewall\\(rule list)};
  \node[box, right=13mm of fw] (svc) {service};
  \draw[flow] (in) -- (fw);
  \draw[flow] (fw) -- (svc) node[midway, above, font=\scriptsize]{allowed :80, :443};
  \draw[-{Stealth[length=2mm]}, kristiania, line width=0.7pt] (fw.south) -- ++(0,-0.9)
    node[below, font=\scriptsize, text=kristiania]{everything else: dropped};
\end{tikzpicture}
\caption{A firewall under default-deny. Only traffic an explicit rule permits
(here, the web ports) reaches the service; the final catch-all rule drops the
rest, so anything unforeseen fails closed.}
\label{fig:firewall}
\end{figure}

\begin{definition}[Firewall, default-deny]\label{def:firewall}
A \textbf{firewall} describes a component that inspects network traffic and
permits or blocks it against a set of rules. Under a \textbf{default-deny}
(allow-list) policy the final, catch-all rule blocks everything, so only
traffic that an earlier rule explicitly permits gets through. The opposite,
default-allow (block-list), permits everything not explicitly forbidden and
fails \emph{open}: anything not foreseen is allowed.
\end{definition}

Default-deny is the stance of the minimum viable network of Session~3: list
the legitimate flows and forbid the rest. A further distinction is
\emph{stateless} versus \emph{stateful}. A stateless firewall judges each
packet alone, whereas a stateful one remembers active connections, so that
having allowed a connection outward, it lets the replies return without a
separate rule. Security groups and the host firewall of the laboratory are
stateful.

The host firewall on Ubuntu is \texttt{ufw} (uncomplicated firewall), a front
end to the kernel's packet filter, and its opening sequence encodes
default-deny directly:
\begin{center}
\ttfamily
sudo ufw default deny incoming\\
sudo ufw default allow outgoing\\
sudo ufw allow 22/tcp\\
sudo ufw allow 80/tcp\\
sudo ufw enable
\end{center}
\noindent The sequence refuses all incoming traffic by default, permits
outgoing traffic, adds two exceptions (SSH and HTTP) and switches the policy
on. Afterwards \texttt{sudo ufw status} prints the loaded policy, and it is
always read back, because the gap between the rules one believes one wrote
and the rules that are loaded is where incidents hide. Three mistakes are
common. The first is leaving SSH open to the whole Internet rather than to
the administrative address, the risk of Example~\ref{ex:ssh-trc}. The second
is opening a database port for testing and forgetting it. The third is
assuming closed when the default is open, as on Azure.

\begin{example}[Example~\ref{ex:ssh-trc}, implemented]\label{ex:ufw-trc}
Session~1 named a threat, judged its risk and prescribed a control for the
open SSH port; here the control is built. \emph{Threat}: automated scanners
brute-force the SSH login within minutes of the port appearing. \emph{Risk}:
likelihood high (constant scanning), impact high (a login is full control).
\emph{Control}: allow SSH only from the administrator's own address, with
\texttt{sudo ufw allow from \emph{YOUR\_IP} to any port 22 proto tcp} instead
of the blanket \texttt{allow 22/tcp} above, and keep the web port open to all.
One \texttt{ufw status} now shows which CIA property each rule protects and
whose responsibility the rule was.
\end{example}

\section{Firewall types, NGFW and WAF}

A plain firewall (\texttt{ufw}, a security group, an Azure NSG) judges
traffic by address, port and protocol, and it cannot see inside an allowed
connection: once port~443 is permitted, everything on port~443 passes,
malicious payloads included. The firewall family is therefore ordered by
how deeply each member looks. The \emph{packet-filter} firewall judges each
packet in isolation by address, port and protocol: stateless, fast, and
blind to context. The \emph{stateful} firewall adds memory of active
connections, and \texttt{ufw} and every cloud security group are of this
kind. Both stop at Layer~4, so they know where traffic goes but not what it
carries. The \emph{next-generation firewall} (NGFW) looks inside: it
performs \emph{deep packet inspection} (Comer, 2021, p.~235) to identify the
application in a flow, applies rules by user identity rather than only by
address, integrates \emph{intrusion prevention} (Session~9), and consults
threat-intelligence feeds of known-bad addresses and domains. Palo Alto,
FortiGate and Cisco Firepower are on-premises examples; Azure Firewall
Premium and AWS Network Firewall are cloud examples.

In the cloud the choice is one of deployment (CSA, 2024, p.~159). A
\emph{CSP firewall} such as Azure Firewall or AWS Network Firewall is built
into the platform, so there is nothing to maintain, but it is limited in
customisation and advanced features. A \emph{virtual appliance}, a vendor
NGFW or IDS/IPS run as a VM, offers more control at the cost of patching and
scaling it oneself. The limit of the managed firewall is concrete: Azure
Firewall's application rules filter by domain name only, without TLS
termination or deep packet inspection, and the vendor documentation itself
calls them very limited compared with a next-generation firewall such as a
Palo Alto appliance.

The \emph{web application firewall} (WAF) is not a network firewall and must
not be confused with one. It sits in front of a single web application and
inspects HTTP requests at Layer~7 for application-layer attacks (SQL
injection, cross-site scripting, the application-layer floods of
Session~3), matched against rule sets such as the OWASP core rules (CSA,
2024, p.~159). An NGFW thus secures the network broadly, whereas a WAF
secures one web application specifically. Table~\ref{tab:firewalls}
compares them.

\begin{table}[htbp]
\centering
\small
\begin{tabular}{@{}l l p{4.3cm} p{3.3cm}@{}}
\toprule
\textbf{Tool} & \textbf{Layer} & \textbf{Inspects} & \textbf{Example} \\
\midrule
Packet-filter firewall & L3--L4 & address, port, protocol; each packet alone (stateless) & early \texttt{iptables} rule \\
Stateful firewall & L3--L4 & the above plus connection state (replies allowed back) & \texttt{ufw}, cloud security group, Azure NSG \\
Next-gen firewall (NGFW) & L3--L7 & application identity, payload (DPI), user, threat intelligence; IPS built in & Palo Alto, Azure Firewall Premium \\
Web app firewall (WAF) & L7 (HTTP) & HTTP request content of one web application (SQLi, XSS, floods) & Azure App Gateway WAF, ModSecurity \\
\bottomrule
\end{tabular}
\caption{Firewalls and the WAF, ordered by inspection depth. A WAF is not a
better firewall: it guards one application at Layer~7, where a network
firewall does not operate.}
\label{tab:firewalls}
\end{table}

Figure~\ref{fig:fwlayers} shows the same comparison as a grid of layers; a
filled cell means the tool operates there.

\begin{figure}[htbp]
\centering
\begin{tikzpicture}[font=\footnotesize]
  \tikzset{on/.style={fill=kristiania!16, draw=kristiania, line width=0.5pt},
           off/.style={fill=klight, draw=kmid, line width=0.4pt}}
  % column headers (layers) -- widely spaced, smaller font, no overlap
  \node[font=\footnotesize\bfseries, text=kgrey, align=center] at (1.9,2.35) {L3\\Network (IP)};
  \node[font=\footnotesize\bfseries, text=kgrey, align=center] at (5.2,2.35) {L4\\Transport (ports)};
  \node[font=\footnotesize\bfseries, text=kgrey, align=center] at (8.5,2.35) {L7\\Application (HTTP)};
  % row: packet filter
  \node[font=\footnotesize, text=kgrey, anchor=east] at (0.55,1.60) {packet filter};
  \filldraw[on]  (0.9,1.35) rectangle (2.9,1.85);
  \filldraw[on]  (4.2,1.35) rectangle (6.2,1.85);
  \filldraw[off] (7.5,1.35) rectangle (9.5,1.85);
  % row: stateful
  \node[font=\footnotesize, text=kgrey, anchor=east] at (0.55,0.95) {stateful};
  \filldraw[on]  (0.9,0.70) rectangle (2.9,1.20);
  \filldraw[on]  (4.2,0.70) rectangle (6.2,1.20);
  \filldraw[off] (7.5,0.70) rectangle (9.5,1.20);
  % row: NGFW
  \node[font=\footnotesize, text=kgrey, anchor=east] at (0.55,0.30) {NGFW};
  \filldraw[on]  (0.9,0.05) rectangle (2.9,0.55);
  \filldraw[on]  (4.2,0.05) rectangle (6.2,0.55);
  \filldraw[on]  (7.5,0.05) rectangle (9.5,0.55);
  % row: WAF
  \node[font=\footnotesize, text=kgrey, anchor=east] at (0.55,-0.35) {WAF};
  \filldraw[off] (0.9,-0.60) rectangle (2.9,-0.10);
  \filldraw[off] (4.2,-0.60) rectangle (6.2,-0.10);
  \filldraw[on]  (7.5,-0.60) rectangle (9.5,-0.10);
\end{tikzpicture}
\caption{Which layers each tool operates at (filled~=~yes). Packet-filter
and stateful firewalls act at Layers~3--4 (addresses, ports); an NGFW reaches
up through Layer~7 with deep packet inspection; a WAF operates only at
Layer~7, on the content of one web application.}
\label{fig:fwlayers}
\end{figure}

One idea links the WAF to the intrusion prevention of Session~9. A WAF runs
either in \emph{detection} mode, where it logs matching requests and lets
them through, or in \emph{prevention} mode, where it blocks them, and that
is the same choice that separates an IDS from an IPS. A WAF prevents attacks
from reaching an application, but it cannot repair a flaw that is already
there, and it is no substitute for writing the application properly (CSA,
2024, p.~247). Session~8 covers the WAF in full; on the students' service it
runs in front of Nginx.

\begin{realworld}{Capital One, 2019}
A chain of misconfigurations exposed about 106~million credit-card
applications. The entry point was a misconfigured web application firewall
running on a cloud instance whose attached identity held far more permission
than it needed. A server-side request forgery flaw made the WAF query the
cloud platform's internal metadata service, which returned the temporary
credentials of that identity. The credentials could list and copy the
contents of more than 700 storage buckets. Three facts matter for this course. A security control can itself become
the attack vector. The damage was set not by the WAF but by the excessive
privilege of the identity behind it, a least-privilege failure (Session~7).
And standard monitoring did not flag the theft, because it looked like
ordinary authorised API calls (Session~9).
\end{realworld}

\section{Defence in depth and Zero Trust}

No single control is sufficient, so controls are layered. \emph{Defence in
depth} means independent layers, so that the failure of one does not expose
everything behind it: network segmentation on the outside, host firewalls
next, then application controls, identities, and encrypted data at the
centre. The tiered minimum viable network of Session~3 is one instance. A
\emph{bastion host}, a single hardened and closely watched machine that is
the only entry point for administration, is another, and segmenting the
network into tiers that may speak only to their neighbours is a third.

\begin{figure}[htbp]
\centering
\begin{tikzpicture}
  \draw[kgrey] (0,0) circle (2.65);
  \draw[kgrey] (0,0) circle (2.05);
  \draw[kgrey] (0,0) circle (1.5);
  \draw[kgrey] (0,0) circle (0.95);
  \node[redbox, inner sep=3pt] at (0,0) {data};
  \node[font=\footnotesize, text=kgrey] at (0,1.22) {identity};
  \node[font=\footnotesize, text=kgrey] at (0,1.77) {application / WAF};
  \node[font=\footnotesize, text=kgrey] at (0,2.32) {host firewall};
  \node[font=\footnotesize, text=kgrey] at (0,2.95) {network segmentation};
\end{tikzpicture}
\caption{Defence in depth. Independent layers around the asset, so that the
failure of any one does not expose everything behind it. Each later session
adds a ring.}
\label{fig:depth}
\end{figure}

\begin{definition}[Zero Trust]\label{def:zerotrust}
\textbf{Zero Trust} describes a security strategy in which nobody and nothing
is trusted merely because of where it sits on the network. It assumes that a
breach has happened or will happen, and it therefore checks every user, device,
application and transaction again and again, no matter whether a request
comes from inside the network or from outside (CSA, 2024, p.~168).
\end{definition}

The Guidance contrasts Zero Trust with the traditional castle-and-moat
perimeter and regards the latter as no longer effective now that services
live in the cloud and staff work remotely (CSA, 2024, p.~50). In the old
model a hard perimeter guarded a trusting interior, so an attacker who
breached the wall could move freely. Comer (2021) gives the cloud reason for
abandoning it (Comer, 2021, pp.~235--237). Traditional security rested on a
firm split between insiders and outsiders, enforced at the perimeter by
firewalls, deep packet inspection and a \emph{demilitarised zone} that
exposed only a few chosen servers. In the cloud, however, the provider's
placement algorithms run many tenants' virtual machines on the same physical
hardware, and even with separate virtual networks their traffic crosses the
same physical links. A tenant therefore has no physical perimeter and cannot
sort the world into inside and outside. Privileges must instead be assigned
per service and every request validated separately, using the identity of
the user and the device, and nobody is trusted for having logged in once.

In the NIST reference model (SP~800-207) a \emph{policy decision point}
decides and a \emph{policy enforcement point} enforces, both placed in the
path of every access (CSA, 2024, p.~171). Supporting technologies include
the software-defined perimeter, which hides resources from the network until
a client has been authenticated (CSA, 2024, p.~172), and, at the network
edge, SASE (CSA, 2024, p.~174). Students are not asked to deploy this stack,
only to recognise its principles: least privilege, verify explicitly, assume
breach.

\begin{example}[The red thread, one week early]\label{ex:mvn-preview}
The VM this week runs one exposed service (Nginx) and one administrative door
(SSH), and the firewall closes everything else. In Session~5 Nextcloud arrives
as a set of tiers (web server, application, database), and the default-deny
stance scales up to the minimum viable network of Session~3. Each tier then
accepts traffic only from the tier above it, and the database is reachable
from nowhere but the application. No new principles are needed for that, only more
rules. Defence in depth is therefore not a product to buy, but default-deny
applied at every layer one owns.
\end{example}

\paragraph{Reading for this chapter.} Comer, Chapter~7 (Virtual Networks:
overlays, VLANs, VXLAN, NAT, SDN and OpenFlow, pp.~87--96) and Chapter~17,
\S\S17.3--17.5 (traditional perimeter security, why it fails in the cloud, and
the Zero Trust model, pp.~235--237); CSA Security Guidance, Domain~7,
Section~7.2 (Cloud Network Fundamentals, including security groups and NACLs,
pp.~153--162) and Section~7.4 (Zero Trust \& Secure Access Service Edge,
pp.~168--176). Optional deepening: J\o sang, Chapter~6 (Network Security, incl.\ firewalls in \S6.4, pp.~117--142). The
Microsoft Learn module is \emph{Describe Azure networking services}; students
map the session's ideas onto Virtual Networks, Network Security Groups and
Azure Firewall.

\section{Summary}

Machines address one another by IP and reach services by port. The pair,
with the protocol, is what a firewall rule matches, and every listening port
is part of the attack surface (Definition~\ref{def:attacksurface}). In the
cloud the network is software: overlays on a shared underlay, programmable
through SDN (Definition~\ref{def:sdn}) and secured in the fabric by stateful
security groups (Definition~\ref{def:secgroup}), whose defaults differ by
provider and are never assumed. The firewall (Definition~\ref{def:firewall})
runs default-deny, and in the laboratory a handful of \texttt{ufw} rules
close the SSH door of Chapter~1 (Example~\ref{ex:ufw-trc}). Plain filters
cannot see inside allowed traffic, so an NGFW adds application awareness and
deep packet inspection, and a WAF guards one web application. Defence in
depth layers these controls, and Zero Trust (Definition~\ref{def:zerotrust})
removes the trusted interior. Session~5 installs Nextcloud on the machine.

\section*{Self-check}
\addcontentsline{toc}{section}{Self-check}

Sketch answers follow the exercises.

\begin{exercise}[Attack surface]
Explain the term attack surface in one's own words, and state why
\texttt{ss -tulpn} is a security command and not merely an administrator's
one.
\end{exercise}

\begin{exercise}[Default-deny reasoning]
State the default-deny principle, and explain why it is the same reasoning as
the minimum viable network of Session~3.
\end{exercise}

\begin{exercise}[A cloud firewall default]
A team member says: ``I added an allow rule on Azure, so now only the web port
is open.'' Why might the service still be exposed, and what does the answer
have to do with AWS behaving differently?
\end{exercise}

\begin{exercise}[NGFW vs WAF]
Distinguish an NGFW from a WAF in one sentence each, and state which one
applies against an application-layer flood on a Nextcloud login page.
\end{exercise}

\begin{exercise}[SSH brute force as threat--risk--control]
Write, in the vocabulary of Definition~\ref{def:trc}, the threat, risk and
control for an SSH port left open to the whole Internet, and give the
\texttt{ufw} rule that implements the control.
\end{exercise}

\paragraph{Sketch answers.}
{\small
(1)~The attack surface is every point an attacker can interact with: open
ports, reachable services, exposed endpoints; \texttt{ss -tulpn} enumerates the
listening ports, so it measures that surface and allows each door to be
justified. (2)~Default-deny blocks everything except explicitly permitted
traffic; like the MVN, it works by listing the few legitimate flows and
forbidding the rest, which needs no knowledge of the attacker. (3)~Azure's
network security groups allow east--west traffic within the virtual network by
default (while denying inbound from the Internet), so intra-VNet paths may be
open even when no rule was added; on AWS a security group denies all inbound
traffic, including between resources in the same group, until it is allowed.
The mental model is the same, namely know the default before trusting it, and
assuming the wrong one can leave a resource reachable from an unintended
location. (4)~An NGFW secures the network broadly with application awareness
and deep packet inspection; a WAF secures one web application by inspecting
HTTP requests, so against an application-layer flood on the login page the WAF
applies (Session~8). (5)~Threat: automated brute-forcing of the SSH login.
Risk: high likelihood (constant scanning), high impact (full control on
success). Control: restrict SSH to the administrative address with
\texttt{sudo ufw allow from \emph{YOUR\_IP} to any port 22 proto tcp}.}


